% GENERATED FILE. Edit canonical lessons, then run npm run generate:books.
% canonical-source-hash: fnv1a64:e5b19a3d9da144d8
% canonical-lessons: TE-C422-hami, TE-C422-neravercu, TE-C422-uddesincu, TE-C422-sankalpam, TE-C422-niscayincu

\chapter{Promises and Intentions}
\label{ch:te-promises-and-intentions}
\hlchaptermodality{422}

\begin{chapteropening}
\textbf{By the end of this chapter:} \emph{I can talk about guarantees, intentions and keeping my word.}

Complete the last lesson of chapter 422: I can talk about guarantees, intentions and keeping my word.
\end{chapteropening}

\section[\texorpdfstring{\te{హామీ} (hāmī)}{హామీ (hāmī)}]{\te{హామీ} (hāmī) — a guarantee; an assurance}

\label{lesson:TE-C422-hami}



\begin{quote}
Before the new one: say the Telugu for to prevent, then the Telugu for a dilemma.
\end{quote}

\subsection*{You'll want to know: \te{హామీ}}
\textbf{\te{హామీ}} — \emph{hāmī} — \textquotedblleft{}a guarantee; an assurance\textquotedblright{}.

\begin{quote}
\te{ప్రభుత్వం} \te{రైతులకు} \te{సహాయం} \te{చేస్తామని} \te{హామీ} \te{ఇచ్చింది}.
\end{quote}

\begin{quote}
\emph{prabhutvaṁ raitulaku sahāyaṁ cēstāmani hāmī iccindi.}
\end{quote}

\begin{quote}
The government assured the farmers that it would help.
\end{quote}

\te{హామీ} \te{ఇచ్చు}, to promise, to guarantee.

\begin{cousinweb}
\textbf{\te{హామీ}} (\emph{hāmī}) has an origin that is not settled; it may have come into Telugu through Hindustani.
\end{cousinweb}

\begin{grammarlens}[title={What you've built}]
One more word for saying what happens if.
\end{grammarlens}

\subsection*{Your turn}
\begin{itemize}
\raggedright
  \item \emph{Say it:} \emph{hāmī}
  \item \emph{Say it:} \emph{hāmī}, once more
  \item \emph{Say it:} \emph{prabhutvaṁ raitulaku sahāyaṁ cēstāmani hāmī iccindi.} — The government assured the farmers that it would help.
  \item \emph{Recall:} say the Telugu for to prevent, then the Telugu for a dilemma, then say \emph{hāmī} again
\end{itemize}

\begin{tcolorbox}[breakable,colback=teal!4,colframe=teal!35!black,title={Before you move on}]
What does \te{హామీ} mean? (\textquotedblleft{}A guarantee; an assurance\textquotedblright{}.) Say it once more.
\end{tcolorbox}
\section[\texorpdfstring{\te{నెరవేర్చు} (neravērcu)}{నెరవేర్చు (neravērcu)}]{\te{నెరవేర్చు} (neravērcu) — to fulfil; to carry out}

\label{lesson:TE-C422-neravercu}



\begin{quote}
Before the new one: say the Telugu for a dilemma, then the Telugu for a guarantee.
\end{quote}

\subsection*{You'll want to know: \te{నెరవేర్చు}}
\textbf{\te{నెరవేర్చు}} — \emph{neravērcu} — \textquotedblleft{}to fulfil; to carry out\textquotedblright{}.

\begin{quote}
\te{అతను} \te{తన} \te{వాగ్దానాన్ని} \te{నెరవేర్చాడు}.
\end{quote}

\begin{quote}
\emph{atanu tana vāgdānānni neravērcāḍu.}
\end{quote}

\begin{quote}
He kept his promise.
\end{quote}

Past (I …): \te{నెరవేర్చాను}, \emph{neravērcānu}.

For promises, wishes and duties: \te{కోరిక} \te{నెరవేర్చు}, to fulfil a wish.

\begin{cousinweb}
\textbf{\te{నెరవేర్చు}} (\emph{neravērcu}) is a Telugu verb whose origin is uncertain; DEDR has no article for it. Its intransitive partner is \te{నెరవేరు}, to come true.
\end{cousinweb}

\begin{grammarlens}[title={What you've built}]
One more word for saying what happens if.
\end{grammarlens}

\subsection*{Your turn}
\begin{itemize}
\raggedright
  \item \emph{Say it:} \emph{neravērcu}
  \item \emph{Say it:} \emph{neravērcu}, once more
  \item \emph{Say it:} \emph{atanu tana vāgdānānni neravērcāḍu.} — He kept his promise.
  \item \emph{Recall:} say the Telugu for a dilemma, then the Telugu for a guarantee, then say \emph{neravērcu} again
\end{itemize}

\begin{tcolorbox}[breakable,colback=teal!4,colframe=teal!35!black,title={Before you move on}]
What does \te{నెరవేర్చు} mean? (\textquotedblleft{}To fulfil; to carry out\textquotedblright{}.) Say it once more.
\end{tcolorbox}
\section[\texorpdfstring{\te{ఉద్దేశించు} (uddēśiñcu)}{ఉద్దేశించు (uddēśiñcu)}]{\te{ఉద్దేశించు} (uddēśiñcu) — to intend; to mean for}

\label{lesson:TE-C422-uddesincu}



\begin{quote}
Before the new one: say the Telugu for a guarantee, then the Telugu for to fulfil.
\end{quote}

\subsection*{You'll want to know: \te{ఉద్దేశించు}}
\textbf{\te{ఉద్దేశించు}} — \emph{uddēśiñcu} — \textquotedblleft{}to intend; to mean for\textquotedblright{}.

\begin{quote}
\te{నేను} \te{మిమ్మల్ని} \te{బాధపెట్టాలని} \te{ఉద్దేశించలేదు}.
\end{quote}

\begin{quote}
\emph{nēnu mimmalni bādhapeṭṭālani uddēśiñcalēdu.}
\end{quote}

\begin{quote}
I did not mean to hurt you.
\end{quote}

Past (I …): \te{ఉద్దేశించాను}, \emph{uddēśiñcānu}.

The verb of \te{ఉద్దేశం}, an intention, which you know.

\begin{cousinweb}
\textbf{\te{ఉద్దేశించు}} (\emph{uddēśiñcu}) is built on Sanskrit \emph{uddeśa}, \textquotedblleft{}the act of pointing to; direction\textquotedblright{} (Monier-Williams), with the Telugu verb ending -\te{ఇంచు}; \te{ఉద్దేశం} is the same word.
\end{cousinweb}

\begin{grammarlens}[title={What you've built}]
One more word for saying what happens if.
\end{grammarlens}

\subsection*{Your turn}
\begin{itemize}
\raggedright
  \item \emph{Say it:} \emph{uddēśiñcu}
  \item \emph{Say it:} \emph{uddēśiñcu}, once more
  \item \emph{Say it:} \emph{nēnu mimmalni bādhapeṭṭālani uddēśiñcalēdu.} — I did not mean to hurt you.
  \item \emph{Recall:} say the Telugu for a guarantee, then the Telugu for to fulfil, then say \emph{uddēśiñcu} again
\end{itemize}

\begin{tcolorbox}[breakable,colback=teal!4,colframe=teal!35!black,title={Before you move on}]
What does \te{ఉద్దేశించు} mean? (\textquotedblleft{}To intend; to mean for\textquotedblright{}.) Say it once more.
\end{tcolorbox}
\section[\texorpdfstring{\te{సంకల్పం} (saṅkalpaṁ)}{సంకల్పం (saṅkalpaṁ)}]{\te{సంకల్పం} (saṅkalpaṁ) — a resolve; a firm intention}

\label{lesson:TE-C422-sankalpam}



\begin{quote}
Before the new one: say the Telugu for to fulfil, then the Telugu for to intend.
\end{quote}

\subsection*{You'll want to know: \te{సంకల్పం}}
\textbf{\te{సంకల్పం}} — \emph{saṅkalpaṁ} — \textquotedblleft{}a resolve; a firm intention\textquotedblright{}.

\begin{quote}
\te{రోజూ} \te{నడవాలని} \te{నా} \te{సంకల్పం}.
\end{quote}

\begin{quote}
\emph{rōjū naḍavālani nā saṅkalpaṁ.}
\end{quote}

\begin{quote}
My resolve is to walk every day.
\end{quote}

\te{సంకల్పించు}, to resolve.

\begin{cousinweb}
\textbf{\te{సంకల్పం}} (\emph{saṅkalpaṁ}) is borrowed from Sanskrit \emph{saṁkalpa}, \textquotedblleft{}will, desire\textquotedblright{} (CDIAL 12819).
\end{cousinweb}

\begin{grammarlens}[title={What you've built}]
One more word for saying what happens if.
\end{grammarlens}

\subsection*{Your turn}
\begin{itemize}
\raggedright
  \item \emph{Say it:} \emph{saṅkalpaṁ}
  \item \emph{Say it:} \emph{saṅkalpaṁ}, once more
  \item \emph{Say it:} \emph{rōjū naḍavālani nā saṅkalpaṁ.} — My resolve is to walk every day.
  \item \emph{Recall:} say the Telugu for to fulfil, then the Telugu for to intend, then say \emph{saṅkalpaṁ} again
\end{itemize}

\begin{tcolorbox}[breakable,colback=teal!4,colframe=teal!35!black,title={Before you move on}]
What does \te{సంకల్పం} mean? (\textquotedblleft{}A resolve; a firm intention\textquotedblright{}.) Say it once more.
\end{tcolorbox}
\section[\texorpdfstring{\te{నిశ్చయించు} (niścayiñcu)}{నిశ్చయించు (niścayiñcu)}]{\te{నిశ్చయించు} (niścayiñcu) — to decide firmly; to make up one’s mind}

\label{lesson:TE-C422-niscayincu}



\begin{quote}
Before the new one: say the Telugu for to intend, then the Telugu for a resolve.
\end{quote}

\subsection*{You'll want to know: \te{నిశ్చయించు}}
\textbf{\te{నిశ్చయించు}} — \emph{niścayiñcu} — \textquotedblleft{}to decide firmly; to make up one’s mind\textquotedblright{}.

\begin{quote}
\te{మేము} \te{ఆదివారం} \te{బయలుదేరాలని} \te{నిశ్చయించాము}.
\end{quote}

\begin{quote}
\emph{mēmu ādivāraṁ bayaludērālani niścayiñcāmu.}
\end{quote}

\begin{quote}
We made up our minds to set off on Sunday.
\end{quote}

Past (I …): \te{నిశ్చయించాను}, \emph{niścayiñcānu}.

The verb of \te{నిశ్చయం}, certainty; \te{నిశ్చయించుకో}, to resolve for oneself.

\begin{cousinweb}
\textbf{\te{నిశ్చయించు}} (\emph{niścayiñcu}) is built on Sanskrit \emph{niścaya}, \textquotedblleft{}inquiry, conviction\textquotedblright{} (CDIAL 7439), with the Telugu verb ending -\te{ఇంచు}.
\end{cousinweb}

\begin{grammarlens}[title={What you've built}]
One more word for saying what happens if.
\end{grammarlens}

\subsection*{Your turn}
\begin{itemize}
\raggedright
  \item \emph{Say it:} \emph{niścayiñcu}
  \item \emph{Say it:} \emph{niścayiñcu}, once more
  \item \emph{Say it:} \emph{mēmu ādivāraṁ bayaludērālani niścayiñcāmu.} — We made up our minds to set off on Sunday.
  \item \emph{Recall:} say the Telugu for to intend, then the Telugu for a resolve, then say \emph{niścayiñcu} again
\end{itemize}

\begin{tcolorbox}[breakable,colback=teal!4,colframe=teal!35!black,title={Before you move on}]
What does \te{నిశ్చయించు} mean? (\textquotedblleft{}To decide firmly; to make up one’s mind\textquotedblright{}.) Say it once more.
\end{tcolorbox}
